\section{Improving the SAEs: Data-Scaled Retraining and a Fixed-L0 KL Fine-tune}\label{sec:improve}

This section reports the GPT-2 improvement phase (R2, R4) and two reference experiments (R5, E2). Everything here is a
\evlabel{validation} finding, except Section~\ref{sec:m5}, which is the single pre-registered \evlabel{test} pass.
\textbf{All improved SAEs are trained by us.} The public \texttt{gpt2-small-res-jb} SAEs~\citep{bloom2024gpt2} serve only as an external reference
(E2). They are never counted as our improvement.

\subsection{R2: data-scaled retraining}\label{sec:r2}

\paragraph{Problem.} Our base SAEs were trained on 101,600 OWT tokens, at width 3,072, and their validation EV was
still rising when training stopped (Instructor Review \#8). It is therefore unclear which of our findings reflect SAEs
in general and which reflect undertraining.

\paragraph{Solution.} We retrained at layers 1, 8 and 11 with a modern recipe and about 1000$\times$ more data:
\begin{itemize}
\item TopK with an AuxK dead-latent loss~\citep{gao2024scaling}, BatchTopK~\citep{bussmann2024batchtopk}, and ReLU with an L1 coefficient
  controlled to L0\,=\,128;
\item width 12,288, plus a TopK-3072 width ablation;
\item 100M fresh OWT tokens per SAE (shard 78, excluding every document in our own sample);
\item encoder initialised to the decoder transpose, input scaling folded into the weights, cosine decay;
\item three seeds each;
\item a frozen-random-decoder TopK as a random baseline~\citep{korznikov2026random}.
\end{itemize}
A modern JumpReLU (STE) did not settle at L0\,=\,128 in the pilot. It was abandoned and disclosed before the full run.

\paragraph{Pipeline.}
\begin{enumerate}
\item Activations are streamed from GPT-2 into a GPU buffer. 39 SAEs were trained in 5.5\,h on one A10G, and GPT-2's
  parameters were unchanged afterwards.
\item Every SAE (old, R2 and res-jb) goes through one evaluation \emph{ladder}. It imports the existing M2 (information), RLM-RP
  (representation), M4 (causal) and splicing (behaviour) code unchanged. This harness reproduces the committed old-SAE
  results (splicing $\leq 2\times10^{-5}$, M4 $\leq 10^{-6}$, logistic gaps $\leq 8\times10^{-4}$).
\item The comparisons and their rule were fixed and hashed before training:
  \begin{itemize}
  \item IMPROVEMENT requires all 9 new/old seed pairs to favour the new SAE \emph{and} a Welch 95\,\% interval
    excluding 0.
  \item Every SAE had OWT L0 within $\pm$10\,\% of 128, so none was excluded.
  \end{itemize}
\end{enumerate}

\paragraph{Results.} Tables~\ref{tab:r2_c1} and~\ref{tab:r2_c3} give before/after values for TopK and for the best old
reconstructer (JumpReLU) against R2 BatchTopK. Table~\ref{tab:r2_counts} summarises all comparisons.
\begin{itemize}
\item \textbf{Reconstruction and behaviour.} Both improve at every layer and for every architecture.
  \begin{itemize}
  \item NMSE falls by 31--74\,\% and spliced $\Delta$CE by 37--82\,\%.
  \item For TopK at L8, $\Delta$CE falls from 0.126 to 0.060 nats, and at L11 from 0.213 to 0.100.
  \item At equal width (A1), data and recipe alone give 23--35\,\% of the NMSE reduction and 33--38\,\% of the
    $\Delta$CE reduction. Width (A2) gives the rest.
  \end{itemize}
\item \textbf{Causal preservation.} It improves for the first time under a pre-registered rule.
  \begin{itemize}
  \item TopK $T_{\text{rec}}/T_{\text{full}}$: 0.62$\to$0.73 (L8) and 0.62$\to$0.69 (L11).
  \item BatchTopK reaches 0.75 at L8.
  \end{itemize}
\item \textbf{Persona representation.} Pair-direction cosine improves everywhere (L8: 0.39$\to$0.49).
\item \textbf{Persona information.} The re-trained persona gap falls by 27--31\,\% at L8/L11 for TopK. The gap
  \emph{relative to PCA at equal OWT error}, however, never improves, and for ReLU it regresses at L1 and L8.
\item \textbf{GoT.} No GoT information metric changes clearly. At L8 the random-decoder baseline has a smaller GoT
  re-fit gap than the trained SAE (0.052 vs 0.081), so this metric does not discriminate there.
\end{itemize}

\paragraph{Inference.}
\begin{itemize}
\item Undertraining explains a large part of the behavioural and causal unfaithfulness we measured. It explains little
  of the information loss: better reconstruction raises the PCA reference as much as the SAE, and SAEs still do not
  beat PCA at equal error.
\item Splice damage remains 1.25--1.61$\times$ that of same-norm noise at L8/L11, compared with 1.45--1.96$\times$
  before. So the remaining error is still structured~\citep{engels2024darkmatter}.
\end{itemize}

% generated by scripts/shared/build_paper_artifacts.py from results/openwebtext/r2_scaled_saes/comparisons.csv
\begin{table}[t]
\centering
\small
\caption{R2 before/after, TopK-3072 (old) $\rightarrow$ TopK-12288 (R2) (VALIDATION; L0 = 128 $\pm$ 10\% on OWT; mean$\pm$sd). Improvement: positive = better. Class: IMPROVEMENT if all 9 new/old seed pairs favour R2 and the Welch interval excludes 0 (pre-registered). GoT is evaluated at layers 8 and 11 only. Source: \texttt{results/\allowbreak{}openwebtext/\allowbreak{}r2\_scaled\_saes/\allowbreak{}comparisons.csv}.}
\label{tab:r2_c1}
\begin{adjustbox}{max width=\linewidth}
\begin{tabular}{llccccl}
\toprule
Layer & Metric & Before (old, 3 seeds) & After (R2, 3 seeds) & $\Delta$ [Welch 95\% CI] & Improvement & Class \\
\midrule
L1 & OWT NMSE $\downarrow$ & 0.0811$\pm$0.0004 & 0.0255$\pm$0.0001 & -0.0556 [-0.0567, -0.0545] & +68.5\% & IMPROVEMENT \\
 & spliced $\Delta$CE $\downarrow$ & 0.0256$\pm$0.0015 & 0.0063$\pm$0.0008 & -0.0193 [-0.0224, -0.0161] & +75.3\% & IMPROVEMENT \\
 & spliced KL $\downarrow$ & 0.0288$\pm$0.0013 & 0.0073$\pm$0.0001 & -0.0215 [-0.0247, -0.0184] & +74.7\% & IMPROVEMENT \\
 & persona re-fit gap $\downarrow$ & 0.083$\pm$0.012 & 0.067$\pm$0.003 & -0.016 [-0.045, +0.012] & +19.8\% & no clear change \\
 & persona SAE$-$PCA gap $\downarrow$ & 0.007$\pm$0.012 & 0.021$\pm$0.003 & +0.014 [-0.015, +0.043] & -- & no clear change \\
 & pair-direction cos $\uparrow$ & 0.373$\pm$0.003 & 0.507$\pm$0.003 & +0.135 [+0.128, +0.141] & +36.1\% & IMPROVEMENT \\
 & $T_{\text{rec}}/T_{\text{full}}$ $\uparrow$ & 0.931$\pm$0.014 & 0.988$\pm$0.005 & +0.057 [+0.025, +0.088] & +6.1\% & IMPROVEMENT \\
\midrule
L8 & OWT NMSE $\downarrow$ & 0.1643$\pm$0.0005 & 0.0959$\pm$0.0001 & -0.0684 [-0.0696, -0.0671] & +41.6\% & IMPROVEMENT \\
 & spliced $\Delta$CE $\downarrow$ & 0.1259$\pm$0.0043 & 0.0603$\pm$0.0018 & -0.0656 [-0.0747, -0.0565] & +52.1\% & IMPROVEMENT \\
 & spliced KL $\downarrow$ & 0.1439$\pm$0.0019 & 0.0680$\pm$0.0002 & -0.0759 [-0.0806, -0.0712] & +52.8\% & IMPROVEMENT \\
 & persona re-fit gap $\downarrow$ & 0.156$\pm$0.017 & 0.108$\pm$0.009 & -0.048 [-0.082, -0.014] & +30.7\% & IMPROVEMENT \\
 & persona SAE$-$PCA gap $\downarrow$ & 0.033$\pm$0.017 & 0.027$\pm$0.009 & -0.006 [-0.041, +0.028] & -- & no clear change \\
 & GoT re-fit gap $\downarrow$ & 0.061$\pm$0.015 & 0.055$\pm$0.022 & -0.006 [-0.051, +0.040] & +9.0\% & no clear change \\
 & pair-direction cos $\uparrow$ & 0.394$\pm$0.003 & 0.488$\pm$0.001 & +0.094 [+0.086, +0.103] & +24.0\% & IMPROVEMENT \\
 & $T_{\text{rec}}/T_{\text{full}}$ $\uparrow$ & 0.625$\pm$0.039 & 0.726$\pm$0.020 & +0.101 [+0.019, +0.182] & +16.1\% & IMPROVEMENT \\
\midrule
L11 & OWT NMSE $\downarrow$ & 0.1891$\pm$0.0005 & 0.1129$\pm$0.0002 & -0.0762 [-0.0773, -0.0751] & +40.3\% & IMPROVEMENT \\
 & spliced $\Delta$CE $\downarrow$ & 0.2129$\pm$0.0018 & 0.0995$\pm$0.0011 & -0.1134 [-0.1171, -0.1097] & +53.3\% & IMPROVEMENT \\
 & spliced KL $\downarrow$ & 0.2177$\pm$0.0015 & 0.1034$\pm$0.0011 & -0.1143 [-0.1175, -0.1112] & +52.5\% & IMPROVEMENT \\
 & persona re-fit gap $\downarrow$ & 0.149$\pm$0.015 & 0.109$\pm$0.008 & -0.040 [-0.071, -0.008] & +26.7\% & IMPROVEMENT \\
 & persona SAE$-$PCA gap $\downarrow$ & 0.083$\pm$0.015 & 0.070$\pm$0.008 & -0.013 [-0.045, +0.019] & -- & no clear change \\
 & GoT re-fit gap $\downarrow$ & 0.110$\pm$0.015 & 0.100$\pm$0.030 & -0.010 [-0.073, +0.053] & +9.0\% & no clear change \\
 & pair-direction cos $\uparrow$ & 0.432$\pm$0.001 & 0.513$\pm$0.001 & +0.080 [+0.077, +0.083] & +18.6\% & IMPROVEMENT \\
 & $T_{\text{rec}}/T_{\text{full}}$ $\uparrow$ & 0.620$\pm$0.004 & 0.692$\pm$0.002 & +0.073 [+0.064, +0.082] & +11.8\% & IMPROVEMENT \\
\bottomrule
\end{tabular}
\end{adjustbox}
\end{table}

% generated by scripts/shared/build_paper_artifacts.py from results/openwebtext/r2_scaled_saes/comparisons.csv
\begin{table}[t]
\centering
\small
\caption{R2 before/after, JumpReLU-3072 (old) $\rightarrow$ BatchTopK-12288 (R2) (VALIDATION; L0 = 128 $\pm$ 10\% on OWT; mean$\pm$sd). Improvement: positive = better. Class: IMPROVEMENT if all 9 new/old seed pairs favour R2 and the Welch interval excludes 0 (pre-registered). GoT is evaluated at layers 8 and 11 only. Source: \texttt{results/\allowbreak{}openwebtext/\allowbreak{}r2\_scaled\_saes/\allowbreak{}comparisons.csv}.}
\label{tab:r2_c3}
\begin{adjustbox}{max width=\linewidth}
\begin{tabular}{llccccl}
\toprule
Layer & Metric & Before (old, 3 seeds) & After (R2, 3 seeds) & $\Delta$ [Welch 95\% CI] & Improvement & Class \\
\midrule
L1 & OWT NMSE $\downarrow$ & 0.0690$\pm$0.0003 & 0.0182$\pm$0.0001 & -0.0509 [-0.0517, -0.0500] & +73.7\% & IMPROVEMENT \\
 & spliced $\Delta$CE $\downarrow$ & 0.0089$\pm$0.0021 & 0.0036$\pm$0.0007 & -0.0053 [-0.0099, -0.0006] & +59.5\% & IMPROVEMENT \\
 & spliced KL $\downarrow$ & 0.0153$\pm$0.0003 & 0.0040$\pm$0.0001 & -0.0114 [-0.0122, -0.0106] & +74.2\% & IMPROVEMENT \\
 & persona re-fit gap $\downarrow$ & 0.097$\pm$0.004 & 0.067$\pm$0.003 & -0.030 [-0.039, -0.020] & +30.7\% & IMPROVEMENT \\
 & persona SAE$-$PCA gap $\downarrow$ & 0.025$\pm$0.007 & 0.030$\pm$0.003 & +0.004 [-0.010, +0.018] & -- & no clear change \\
 & pair-direction cos $\uparrow$ & 0.384$\pm$0.001 & 0.526$\pm$0.003 & +0.142 [+0.137, +0.148] & +37.1\% & IMPROVEMENT \\
 & $T_{\text{rec}}/T_{\text{full}}$ $\uparrow$ & 0.968$\pm$0.011 & 0.994$\pm$0.002 & +0.026 [+0.001, +0.050] & +2.6\% & IMPROVEMENT \\
\midrule
L8 & OWT NMSE $\downarrow$ & 0.1598$\pm$0.0002 & 0.0907$\pm$0.0001 & -0.0691 [-0.0695, -0.0687] & +43.3\% & IMPROVEMENT \\
 & spliced $\Delta$CE $\downarrow$ & 0.1173$\pm$0.0029 & 0.0525$\pm$0.0026 & -0.0648 [-0.0711, -0.0585] & +55.2\% & IMPROVEMENT \\
 & spliced KL $\downarrow$ & 0.1360$\pm$0.0020 & 0.0605$\pm$0.0003 & -0.0755 [-0.0803, -0.0706] & +55.5\% & IMPROVEMENT \\
 & persona re-fit gap $\downarrow$ & 0.162$\pm$0.004 & 0.126$\pm$0.007 & -0.037 [-0.051, -0.022] & +22.5\% & IMPROVEMENT \\
 & persona SAE$-$PCA gap $\downarrow$ & 0.040$\pm$0.006 & 0.053$\pm$0.007 & +0.014 [-0.002, +0.029] & -- & no clear change \\
 & GoT re-fit gap $\downarrow$ & 0.085$\pm$0.007 & 0.073$\pm$0.015 & -0.012 [-0.044, +0.020] & +14.4\% & no clear change \\
 & pair-direction cos $\uparrow$ & 0.436$\pm$0.002 & 0.504$\pm$0.003 & +0.068 [+0.062, +0.075] & +15.7\% & IMPROVEMENT \\
 & $T_{\text{rec}}/T_{\text{full}}$ $\uparrow$ & 0.665$\pm$0.011 & 0.754$\pm$0.030 & +0.090 [+0.025, +0.154] & +13.5\% & IMPROVEMENT \\
\midrule
L11 & OWT NMSE $\downarrow$ & 0.1897$\pm$0.0002 & 0.1087$\pm$0.0001 & -0.0810 [-0.0814, -0.0806] & +42.7\% & IMPROVEMENT \\
 & spliced $\Delta$CE $\downarrow$ & 0.2246$\pm$0.0069 & 0.0933$\pm$0.0040 & -0.1313 [-0.1454, -0.1172] & +58.5\% & IMPROVEMENT \\
 & spliced KL $\downarrow$ & 0.2317$\pm$0.0021 & 0.0977$\pm$0.0008 & -0.1340 [-0.1386, -0.1294] & +57.8\% & IMPROVEMENT \\
 & persona re-fit gap $\downarrow$ & 0.118$\pm$0.018 & 0.101$\pm$0.008 & -0.018 [-0.055, +0.019] & +15.0\% & no clear change \\
 & persona SAE$-$PCA gap $\downarrow$ & 0.052$\pm$0.018 & 0.067$\pm$0.008 & +0.014 [-0.023, +0.052] & -- & no clear change \\
 & GoT re-fit gap $\downarrow$ & 0.126$\pm$0.003 & 0.099$\pm$0.019 & -0.026 [-0.072, +0.019] & +21.0\% & no clear change \\
 & pair-direction cos $\uparrow$ & 0.478$\pm$0.005 & 0.517$\pm$0.003 & +0.039 [+0.030, +0.049] & +8.2\% & IMPROVEMENT \\
 & $T_{\text{rec}}/T_{\text{full}}$ $\uparrow$ & 0.678$\pm$0.007 & 0.701$\pm$0.006 & +0.024 [+0.008, +0.039] & +3.5\% & IMPROVEMENT \\
\bottomrule
\end{tabular}
\end{adjustbox}
\end{table}

% generated by scripts/shared/build_paper_artifacts.py from results/openwebtext/r2_scaled_saes/comparisons.csv, results/openwebtext/r2_scaled_saes/criteria.json
\begin{table}[t]
\centering
\small
\caption{R2 comparisons over all ladder metrics at layers 1, 8, 11 (VALIDATION). Reductions are the range over layers; the cell counts cover all 17 metrics $\times$ 3 layers (51 cells; GoT at 8 and 11 only). Source: \texttt{results/\allowbreak{}openwebtext/\allowbreak{}r2\_scaled\_saes/\allowbreak{}comparisons.csv}; \texttt{results/\allowbreak{}openwebtext/\allowbreak{}r2\_scaled\_saes/\allowbreak{}criteria.json}.}
\label{tab:r2_counts}
\begin{adjustbox}{max width=\linewidth}
\begin{tabular}{llccccc}
\toprule
ID & Comparison & NMSE reduction & $\Delta$CE reduction & \multicolumn{3}{c}{cells: improved / no change / regressed} \\
\midrule
C1 & TopK: old-3072 $\rightarrow$ R2-12288 & 40--69\% & 52--75\% & 37 & 14 & 0 \\
C2 & ReLU: old-3072 $\rightarrow$ R2-12288 & 31--65\% & 37--82\% & 30 & 18 & 3 \\
C3 & old JumpReLU-3072 $\rightarrow$ R2 BatchTopK-12288 & 43--74\% & 55--60\% & 36 & 15 & 0 \\
A1 & data + recipe, equal width (TopK-3072) & 23--35\% & 33--38\% & 34 & 17 & 0 \\
A2 & width (R2 TopK 3072 $\rightarrow$ 12288) & 20--52\% & 28--60\% & 27 & 21 & 3 \\
\bottomrule
\end{tabular}
\end{adjustbox}
\end{table}


\subsection{R4: KL fine-tuning at fixed L0}\label{sec:r4}

\paragraph{Problem.} IMP-L's output-KL fine-tune improved behaviour, but for ReLU and JumpReLU only with L0 drift
(not creditable). For TopK it showed no effect, on weak SAEs.

\paragraph{Solution.} We fine-tune the strong R2 TopK-12288 SAEs, whose L0 is fixed at exactly 128.
\begin{itemize}
\item The objective is NMSE plus $\alpha\,$KL$(p_{\text{orig}}\Vert p_{\text{spliced}})$, after \citet{karvonen2025kl}.
\item The run is short: 5M tokens, about 5\,\% of the training budget.
\item An MSE-only control continues on identical tokens.
\end{itemize}

\paragraph{Pipeline.}
\begin{itemize}
\item Layers 8 and 11, 3 seeds, learning rate $2\times10^{-5}$.
\item $\alpha$ is set by first-batch scale matching.
\item The same ladder evaluates every SAE.
\item The IMP-L rule applies: better on all 3 seeds and $|\bar d|>2\,$sd, with an EV guard.
\end{itemize}

\paragraph{Results} (Table~\ref{tab:r4}).
\begin{itemize}
\item Spliced $\Delta$CE falls by 18\,\% against the MSE control at both layers (L8: 0.060$\to$0.049; L11:
  0.100$\to$0.081), with KL down 13--17\,\%. The cost is $-0.002$ EV, so the guard passes.
\item Information gaps do not change.
\item Pair-direction cosine drops by 0.012--0.013.
\item At L11, causal preservation \emph{decreases}: 0.694$\to$0.671, \evlabel{regression}, all seeds.
\end{itemize}

\paragraph{Inference.} A behavioural objective improves behaviour creditably once sparsity is fixed. Its gain does not
transfer to information, representation or the causal persona effect, and partly reverses the causal effect. This is a
further dissociation between faithfulness levels, so the KL fine-tune is a behaviour-only tool.

% generated by scripts/shared/build_paper_artifacts.py from results/openwebtext/r4_kl_finetune/comparisons.csv
\begin{table}[t]
\centering
\small
\caption{R4: short KL+MSE fine-tune of the R2 TopK-12288 SAEs at fixed L0 = 128 (5M tokens, 3 seeds; VALIDATION), against an MSE-only control on identical tokens (K1--K3) and against the R2 base (K4). Rule: better on all 3 seeds and $|\bar d| > 2\,$sd; guard EV(KL) $\geq$ EV(MSE)$-0.02$ (passed: $-0.002$ at both layers). Source: \texttt{results/\allowbreak{}openwebtext/\allowbreak{}r4\_kl\_finetune/\allowbreak{}comparisons.csv}.}
\label{tab:r4}
\begin{adjustbox}{max width=\linewidth}
\begin{tabular}{llllcccl}
\toprule
Layer & ID & Metric & Reference & KL arm & paired $\Delta\pm$sd & Improvement & Class \\
\midrule
L8 & K1 & spliced $\Delta$CE $\downarrow$ & MSE-only 0.0601 & 0.0493 & -0.0107$\pm$0.0011 & +17.9\% & PRELIM.\ IMPROVEMENT \\
L8 & K1 & spliced KL $\downarrow$ & MSE-only 0.0677 & 0.0587 & -0.0089$\pm$0.0002 & +13.2\% & PRELIM.\ IMPROVEMENT \\
L8 & K2 & persona re-fit gap $\downarrow$ & MSE-only 0.1088 & 0.1061 & -0.0026$\pm$0.0035 & +2.4\% & no improvement \\
L8 & K2 & GoT re-fit gap $\downarrow$ & MSE-only 0.0485 & 0.0548 & +0.0063$\pm$0.0034 & -12.9\% & no improvement \\
L8 & K2 & persona frozen gap $\downarrow$ & MSE-only 0.1137 & 0.1135 & -0.0002$\pm$0.0020 & +0.2\% & no improvement \\
L8 & K2 & GoT frozen gap $\downarrow$ & MSE-only 0.2404 & 0.2372 & -0.0032$\pm$0.0168 & +1.3\% & no improvement \\
L8 & K3 & $T_{\text{rec}}/T_{\text{full}}$ $\uparrow$ & MSE-only 0.7202 & 0.7096 & -0.0107$\pm$0.0162 & -1.5\% & no improvement \\
L8 & K4 & spliced $\Delta$CE $\downarrow$ & R2 base 0.0603 & 0.0493 & -0.0110$\pm$0.0007 & +18.2\% & PRELIM.\ IMPROVEMENT \\
L8 & K4 & spliced KL $\downarrow$ & R2 base 0.0680 & 0.0587 & -0.0092$\pm$0.0002 & +13.6\% & PRELIM.\ IMPROVEMENT \\
\midrule
L11 & K1 & spliced $\Delta$CE $\downarrow$ & MSE-only 0.0997 & 0.0814 & -0.0183$\pm$0.0012 & +18.3\% & PRELIM.\ IMPROVEMENT \\
L11 & K1 & spliced KL $\downarrow$ & MSE-only 0.1030 & 0.0858 & -0.0172$\pm$0.0005 & +16.7\% & PRELIM.\ IMPROVEMENT \\
L11 & K2 & persona re-fit gap $\downarrow$ & MSE-only 0.1070 & 0.1106 & +0.0035$\pm$0.0071 & -3.3\% & no improvement \\
L11 & K2 & GoT re-fit gap $\downarrow$ & MSE-only 0.0936 & 0.1017 & +0.0081$\pm$0.0103 & -8.7\% & no improvement \\
L11 & K2 & persona frozen gap $\downarrow$ & MSE-only 0.1192 & 0.1215 & +0.0023$\pm$0.0031 & -1.9\% & no improvement \\
L11 & K2 & GoT frozen gap $\downarrow$ & MSE-only 0.2462 & 0.2478 & +0.0016$\pm$0.0079 & -0.6\% & no improvement \\
L11 & K3 & $T_{\text{rec}}/T_{\text{full}}$ $\uparrow$ & MSE-only 0.6943 & 0.6708 & -0.0235$\pm$0.0015 & -3.4\% & \textbf{REGRESSION} \\
L11 & K4 & spliced $\Delta$CE $\downarrow$ & R2 base 0.0995 & 0.0814 & -0.0181$\pm$0.0012 & +18.2\% & PRELIM.\ IMPROVEMENT \\
L11 & K4 & spliced KL $\downarrow$ & R2 base 0.1034 & 0.0858 & -0.0176$\pm$0.0006 & +17.0\% & PRELIM.\ IMPROVEMENT \\
\bottomrule
\end{tabular}
\end{adjustbox}
\end{table}


\subsection{R5: how much of the persona effect can a low-rank intervention carry?}\label{sec:r5}

\paragraph{Problem.} M4's shares (62--68\,\% for the reconstruction, 2--74\,\% for top-$k$ features) had no
reference point.

\paragraph{Solution.} At each M4 site and rank $k$ we compare:
\begin{itemize}
\item a supervised subspace learned on TRAIN pairs (DAS-style; \citealp{chaudhary2024evaluating});
\item PCA of the TRAIN differences;
\item the mass-mean direction;
\item 20 random subspaces;
\item M4's SAE feature patches.
\end{itemize}

\paragraph{Results} (Table~\ref{tab:r5}).
\begin{itemize}
\item The effect is low-rank:
  \begin{itemize}
  \item one learned direction carries 0.92 (L1), 0.64 (L8) and 0.40 (L11);
  \item four learned directions carry 1.02, 0.91 and 0.62.
  \end{itemize}
\item Four SAE features carry 0.10--0.16.
\item At L8 a rank-4 subspace carries more of the effect than the whole SAE reconstruction (0.91 vs 0.62--0.67).
\item Random subspaces carry $\leq$0.13 even at $k=64$.
\end{itemize}

\paragraph{Inference.} The reconstruction's partial share at L8/L11 is a property of the SAE, not of the task. The SAE
basis is far from the best basis at matched rank, although well above random.

% generated by scripts/shared/build_paper_artifacts.py from results/repe/r5_causal_subspace_skyline/summary.csv
\begin{table}[t]
\centering
\small
\caption{R5: share of the persona effect $T_k/T_{\text{full}}$ carried by rank-$k$ subspaces at the M4 sites (VALIDATION; fitted on TRAIN pairs). For reference, the SAE whole reconstruction carries 0.94--0.97 (L1), 0.62--0.67 (L8) and 0.62--0.68 (L11) for the old SAEs (M4). DAS shares above 1 are possible because the single-site full patch does not fully restore $p_H$. Source: \texttt{results/\allowbreak{}repe/\allowbreak{}r5\_causal\_subspace\_skyline/\allowbreak{}summary.csv}.}
\label{tab:r5}
\begin{adjustbox}{max width=\linewidth}
\begin{tabular}{llcccc}
\toprule
Site & $k$ & DAS (learned, 3 seeds) & PCA of TRAIN $\Delta$ & random mean (q95) & SAE top-$k$ features (range over archs) \\
\midrule
L1 slot & 1 & 0.915$\pm$0.008 & 0.430 & -0.000 (0.003) & 0.038--0.073 \\
 & 4 & 1.016$\pm$0.002 & 1.000 & 0.004 (0.020) & 0.114--0.158 \\
 & 16 & 1.013$\pm$0.000 & 1.000 & 0.012 (0.057) & 0.282--0.332 \\
 & 64 & 1.016$\pm$0.003 & 1.000 & 0.041 (0.082) & 0.498--0.597 \\
\midrule
L8 final & 1 & 0.641$\pm$0.005 & 0.358 & 0.002 (0.005) & 0.037--0.054 \\
 & 4 & 0.913$\pm$0.014 & 0.518 & 0.008 (0.014) & 0.095--0.155 \\
 & 16 & 1.116$\pm$0.004 & 0.678 & 0.027 (0.042) & 0.206--0.315 \\
 & 64 & 1.289$\pm$0.002 & 0.826 & 0.110 (0.139) & 0.338--0.477 \\
\midrule
L11 final & 1 & 0.401$\pm$0.007 & 0.379 & 0.002 (0.004) & 0.026--0.054 \\
 & 4 & 0.619$\pm$0.005 & 0.587 & 0.008 (0.015) & 0.079--0.136 \\
 & 16 & 0.759$\pm$0.002 & 0.706 & 0.033 (0.040) & 0.186--0.288 \\
 & 64 & 0.888$\pm$0.002 & 0.839 & 0.127 (0.143) & 0.320--0.461 \\
\bottomrule
\end{tabular}
\end{adjustbox}
\end{table}


\subsection{E2: do our conclusions hold for an external SAE?}\label{sec:e2}

\begin{itemize}
\item \textbf{Conclusions that hold for res-jb} (24,576 latents, 300M tokens; Table~\ref{tab:e2}):
  \begin{itemize}
  \item the frozen GoT loss is mostly probe fragility (H1 holds for GoT);
  \item the persona loss is real (H1 fails for persona, and H2 holds);
  \item the SAE never beats PCA (H3 holds at all four layers);
  \item the reconstruction carries only part of the persona effect at L8 and L11 (H4).
  \end{itemize}
\item \textbf{The one exception} is the GoT re-fit gap at L8, which is small with a CI touching zero, as for our own
  SAEs.
\item \textbf{Inference.} Instructor Review \#8's concern that our SAEs are small and unconverged does not explain
  these conclusions.
\end{itemize}

% generated by scripts/shared/build_paper_artifacts.py from results/shared/e2_public_sae_reference/hypotheses.csv
\begin{table}[t]
\centering
\small
\caption{E2: our headline conclusions tested on the public \texttt{gpt2-small-res-jb} SAEs (ReLU, 24,576 latents, 300M tokens; VALIDATION). H1: frozen loss mostly probe fragility (noise gap $\geq 0.5\times$ SAE gap); H2: re-trained loss real (CI low $>0$, BH); H3: SAE does not beat PCA (SAE gap $\geq$ PCA gap $-0.02$); H4: reconstruction carries only part of the persona effect (CI high $<0.80$, above controls). $^\dagger$ labelled extension layers. Source: \texttt{results/\allowbreak{}shared/\allowbreak{}e2\_public\_sae\_reference/\allowbreak{}hypotheses.csv}.}
\label{tab:e2}
\begin{adjustbox}{max width=\linewidth}
\begin{tabular}{lllll}
\toprule
Hypothesis & Layer & Task & Value & Result \\
\midrule
E2-H1 & L1$^\dagger$ & RepE-LM persona & SAE 0.150 / noise 0.028 & does not hold \\
E2-H1 & L8 & GoT cities & SAE 0.309 / noise 0.247 & HOLDS \\
E2-H1 & L8 & RepE-LM persona & SAE 0.247 / noise 0.075 & does not hold \\
E2-H1 & L10 & GoT cities & SAE 0.308 / noise 0.242 & HOLDS \\
E2-H1 & L10 & RepE-LM persona & SAE 0.242 / noise 0.058 & does not hold \\
E2-H1 & L11$^\dagger$ & GoT cities & SAE 0.312 / noise 0.222 & HOLDS \\
E2-H1 & L11$^\dagger$ & RepE-LM persona & SAE 0.250 / noise 0.062 & does not hold \\
\midrule
E2-H2 & L1$^\dagger$ & RepE-LM persona & 0.134 (CI low +0.109, $q$ 0.001) & HOLDS \\
E2-H2 & L8 & GoT cities & 0.025 (CI low -0.001, $q$ 0.031) & does not hold \\
E2-H2 & L8 & RepE-LM persona & 0.223 (CI low +0.190, $q$ 0.001) & HOLDS \\
E2-H2 & L10 & GoT cities & 0.132 (CI low +0.102, $q$ 0.001) & HOLDS \\
E2-H2 & L10 & RepE-LM persona & 0.224 (CI low +0.188, $q$ 0.001) & HOLDS \\
E2-H2 & L11$^\dagger$ & GoT cities & 0.144 (CI low +0.110, $q$ 0.001) & HOLDS \\
E2-H2 & L11$^\dagger$ & RepE-LM persona & 0.234 (CI low +0.194, $q$ 0.001) & HOLDS \\
\midrule
E2-H3 & L1$^\dagger$ & RepE-LM persona & SAE 0.134 / PCA 0.068 & HOLDS \\
E2-H3 & L8 & GoT cities & SAE 0.025 / PCA 0.012 & HOLDS \\
E2-H3 & L8 & RepE-LM persona & SAE 0.223 / PCA 0.112 & HOLDS \\
E2-H3 & L10 & GoT cities & SAE 0.132 / PCA 0.028 & HOLDS \\
E2-H3 & L10 & RepE-LM persona & SAE 0.224 / PCA 0.067 & HOLDS \\
E2-H3 & L11$^\dagger$ & GoT cities & SAE 0.144 / PCA 0.042 & HOLDS \\
E2-H3 & L11$^\dagger$ & RepE-LM persona & SAE 0.234 / PCA 0.066 & HOLDS \\
\midrule
E2-H4 & L1$^\dagger$ & RepE-LM persona (M4 site) & 0.979 (CI high 0.987) & does not hold \\
E2-H4 & L8 & RepE-LM persona (M4 site) & 0.559 (CI high 0.601) & HOLDS \\
E2-H4 & L11$^\dagger$ & RepE-LM persona (M4 site) & 0.612 (CI high 0.644) & HOLDS \\
\bottomrule
\end{tabular}
\end{adjustbox}
\end{table}


\subsection{M5: confirmatory test pass}\label{sec:m5}

\paragraph{Protocol.} Before any test row was read, we froze and hashed the following, and then ran it once:
\begin{itemize}
\item the script and every module it imports;
\item the full SAE list: 81 SAEs, with no selection;
\item the metrics;
\item the rules (C1--C3, A1, A2: the R2 rule; K1, K4: the IMP-L rule).
\end{itemize}
It reads OWT test (12,700 tokens per layer) and GoT cities+neg\_cities test (452 rows): 552 test rows in total.
RepE-LM has no test split, so the persona, representation and causal results stay \evlabel{validation}. A smoke run
on validation rows reproduced the ladder exactly for the OWT metrics and to within $10^{-4}$ for GoT.

\paragraph{Results} (\evlabel{test}; Tables~\ref{tab:m5} and~\ref{tab:m5_counts}).
\begin{itemize}
\item \textbf{R2: every reconstruction and behaviour cell replicates.}
  \begin{itemize}
  \item NMSE falls by 31--74\,\% and spliced $\Delta$CE by 35--79\,\% at layers 1, 8 and 11, for all three
    comparisons and both ablations.
  \item For TopK, L8 $\Delta$CE falls from 0.133 to 0.062 nats and L11 from 0.204 to 0.096.
  \end{itemize}
\item \textbf{GoT information shows no confirmed improvement.}
  \begin{itemize}
  \item No comparison improves the GoT re-fit gap or the SAE$-$PCA gap.
  \item Width alone (A2) makes both slightly \emph{worse} at L8 (\evlabel{regression}).
  \end{itemize}
\item \textbf{R4:} the KL fine-tune lowers $\Delta$CE by 16--17\,\% against the MSE control at both layers, at a
  2\,\% NMSE cost (classified as a regression in NMSE, and within the pre-registered EV guard). It has no GoT effect.
\end{itemize}

\paragraph{Inference.} The improvement phase's behavioural and reconstruction gains are \evlabel{confirmed}. The
absence of a truth-information gain is also confirmed: the SAEs reconstruct and behave much better but still do not
preserve GoT truth information better than PCA at equal error.

% generated by scripts/shared/build_paper_artifacts.py from results/shared/m5_confirmatory_test/comparisons.csv
\begin{table}[t]
\centering
\small
\caption{M5: the single pre-registered confirmatory TEST pass (OWT test, 100 sequences, 12,700 tokens; GoT cities+neg\_cities test). Same rules as on validation (C/A: all 9 seed pairs and Welch CI; K: all 3 paired seeds and $|\bar d|>2\,$sd). Means over 3 seeds. C1: TopK old-3072 $\rightarrow$ R2-12288; C3: old JumpReLU-3072 $\rightarrow$ R2 BatchTopK-12288; K1: R4 KL fine-tune vs MSE control. All comparisons are counted in Table~\ref{tab:m5_counts}. Source: \texttt{results/\allowbreak{}shared/\allowbreak{}m5\_confirmatory\_test/\allowbreak{}comparisons.csv}.}
\label{tab:m5}
\begin{adjustbox}{max width=\linewidth}
\begin{tabular}{lllcccl}
\toprule
ID & Layer & Metric & Before & After & Improvement & Class (TEST) \\
\midrule
C1 & L1 & OWT NMSE $\downarrow$ & 0.0787 & 0.0248 & +68.6\% & IMPROVEMENT \\
 & L1 & spliced $\Delta$CE $\downarrow$ & 0.0221 & 0.0056 & +74.9\% & IMPROVEMENT \\
 & L8 & OWT NMSE $\downarrow$ & 0.1598 & 0.0941 & +41.1\% & IMPROVEMENT \\
 & L8 & spliced $\Delta$CE $\downarrow$ & 0.1330 & 0.0620 & +53.4\% & IMPROVEMENT \\
 & L8 & GoT re-fit gap $\downarrow$ & 0.0607 & 0.0558 & +8.1\% & no clear change \\
 & L11 & OWT NMSE $\downarrow$ & 0.1848 & 0.1112 & +39.8\% & IMPROVEMENT \\
 & L11 & spliced $\Delta$CE $\downarrow$ & 0.2041 & 0.0962 & +52.9\% & IMPROVEMENT \\
 & L11 & GoT re-fit gap $\downarrow$ & 0.1009 & 0.0775 & +23.2\% & no clear change \\
\midrule
C3 & L1 & OWT NMSE $\downarrow$ & 0.0675 & 0.0178 & +73.7\% & IMPROVEMENT \\
 & L1 & spliced $\Delta$CE $\downarrow$ & 0.0088 & 0.0019 & +78.8\% & IMPROVEMENT \\
 & L8 & OWT NMSE $\downarrow$ & 0.1560 & 0.0897 & +42.5\% & IMPROVEMENT \\
 & L8 & spliced $\Delta$CE $\downarrow$ & 0.1212 & 0.0542 & +55.3\% & IMPROVEMENT \\
 & L8 & GoT re-fit gap $\downarrow$ & 0.0584 & 0.0654 & -12.0\% & no clear change \\
 & L11 & OWT NMSE $\downarrow$ & 0.1851 & 0.1073 & +42.0\% & IMPROVEMENT \\
 & L11 & spliced $\Delta$CE $\downarrow$ & 0.2155 & 0.0927 & +57.0\% & IMPROVEMENT \\
 & L11 & GoT re-fit gap $\downarrow$ & 0.1070 & 0.0871 & +18.6\% & no clear change \\
\midrule
K1 & L8 & OWT NMSE $\downarrow$ & 0.0938 & 0.0958 & -2.1\% & \textbf{REGRESSION} \\
 & L8 & spliced $\Delta$CE $\downarrow$ & 0.0623 & 0.0514 & +17.4\% & IMPROVEMENT \\
 & L8 & GoT re-fit gap $\downarrow$ & 0.0585 & 0.0484 & +17.3\% & no clear change \\
 & L11 & OWT NMSE $\downarrow$ & 0.1109 & 0.1133 & -2.2\% & \textbf{REGRESSION} \\
 & L11 & spliced $\Delta$CE $\downarrow$ & 0.0964 & 0.0807 & +16.3\% & IMPROVEMENT \\
 & L11 & GoT re-fit gap $\downarrow$ & 0.0780 & 0.0784 & -0.5\% & no clear change \\
\bottomrule
\end{tabular}
\end{adjustbox}
\end{table}

% generated by scripts/shared/build_paper_artifacts.py from results/shared/m5_confirmatory_test/comparisons.csv
\begin{table}[t]
\centering
\small
\caption{M5 (TEST): pre-registered classes over all OWT and GoT cells (layers 1, 8, 11; GoT at 8 and 11). Improvement ranges are over layers; positive = better. Source: \texttt{results/\allowbreak{}shared/\allowbreak{}m5\_confirmatory\_test/\allowbreak{}comparisons.csv}.}
\label{tab:m5_counts}
\begin{adjustbox}{max width=\linewidth}
\begin{tabular}{llcccp{0.25\linewidth}}
\toprule
ID & Comparison & NMSE improvement & $\Delta$CE improvement & improved / no change / regressed & regressed cells \\
\midrule
C1 & TopK: old-3072 $\rightarrow$ R2-12288 & +40 to +69\% & +53 to +75\% & 17 / 7 / 0 &  \\
C2 & ReLU: old-3072 $\rightarrow$ R2-12288 & +31 to +64\% & +35 to +73\% & 17 / 7 / 0 &  \\
C3 & old JumpReLU-3072 $\rightarrow$ R2 BatchTopK-12288 & +42 to +74\% & +55 to +79\% & 17 / 7 / 0 &  \\
A1 & data + recipe, equal width (TopK-3072) & +22 to +35\% & +34 to +40\% & 17 / 7 / 0 &  \\
A2 & width (R2 TopK 3072 $\rightarrow$ 12288) & +20 to +52\% & +26 to +58\% & 16 / 6 / 2 & GoT refit logistic gap L8; GoT sae minus pca refit gap L8 \\
K1 & R4 KL fine-tune vs MSE control & -2 to -2\% & +16 to +17\% & 4 / 2 / 2 & OWT NMSE L11; OWT NMSE L8 \\
K4 & R4 KL fine-tune vs R2 base & -2 to -2\% & +16 to +17\% & 4 / 2 / 2 & OWT NMSE L11; OWT NMSE L8 \\
\bottomrule
\end{tabular}
\end{adjustbox}
\end{table}

